% GENERATED FILE. Edit canonical lessons, then run npm run generate:books.
% canonical-source-hash: fnv1a64:769e6954842fdee9
% canonical-lessons: MR-C505-kumbhar, MR-C505-lohar, MR-C505-sonar, MR-C505-vinkar, MR-C505-thatne

\chapter{Old Crafts}
\label{ch:mr-old-crafts}
\hlchaptermodality{505}

\begin{chapteropening}
\textbf{By the end of this chapter:} \emph{I can name some old crafts and the people who practise them.}

Complete the last lesson of chapter 505: I can name some old crafts and the people who practise them.
\end{chapteropening}

\section[\texorpdfstring{\mr{कुंभार} (kumbhār)}{कुंभार (kumbhār)}]{\mr{कुंभार} (kumbhār) — a potter}

\label{lesson:MR-C505-kumbhar}



\begin{quote}
Before the new one: say the Marathi for a barber, then the Marathi for to hammer in.
\end{quote}

\subsection*{You'll want to know: \mr{कुंभार}}
\textbf{\mr{कुंभार}} — \emph{kumbhār} — \textquotedblleft{}a potter\textquotedblright{}.

\begin{quote}
\mr{कुंभार} \mr{मातीची} \mr{भांडी} \mr{बनवतो}.
\end{quote}

\begin{quote}
\emph{kumbhār mātīcī bhāṇḍī banavto.}
\end{quote}

\begin{quote}
The potter makes clay pots.
\end{quote}

Masculine.

\begin{cousinweb}
\textbf{\mr{कुंभार}} (\emph{kumbhār}) is a tadbhava, inherited from Sanskrit \emph{kumbhakāra}, \textquotedblleft{}potter\textquotedblright{} (CDIAL 3310), from \emph{kumbha}, \textquotedblleft{}a pot\textquotedblright{}, and \emph{kāra}, \textquotedblleft{}maker\textquotedblright{}.
\end{cousinweb}

\begin{grammarlens}[title={What you've built}]
One more word for describing what you lived and felt.
\end{grammarlens}

\subsection*{Your turn}
\begin{itemize}
\raggedright
  \item \emph{Say it:} \emph{kumbhār}
  \item \emph{Say it:} \emph{kumbhār}, once more
  \item \emph{Say it:} \emph{kumbhār mātīcī bhāṇḍī banavto.} — The potter makes clay pots.
  \item \emph{Recall:} say the Marathi for a barber, then the Marathi for to hammer in, then say \emph{kumbhār} again
\end{itemize}

\begin{tcolorbox}[breakable,colback=teal!4,colframe=teal!35!black,title={Before you move on}]
What does \mr{कुंभार} mean? (\textquotedblleft{}A potter\textquotedblright{}.) Say it once more.
\end{tcolorbox}
\section[\texorpdfstring{\mr{लोहार} (lohār)}{लोहार (lohār)}]{\mr{लोहार} (lohār) — a blacksmith}

\label{lesson:MR-C505-lohar}



\begin{quote}
Before the new one: say the Marathi for to hammer in, then the Marathi for a potter.
\end{quote}

\subsection*{You'll want to know: \mr{लोहार}}
\textbf{\mr{लोहार}} — \emph{lohār} — \textquotedblleft{}a blacksmith\textquotedblright{}.

\begin{quote}
\mr{लोहाराने} \mr{शेतासाठी} \mr{नवीन} \mr{अवजार} \mr{बनवलं}.
\end{quote}

\begin{quote}
\emph{lohārāne śetāsāṭhī navīn avjār banavlaṁ.}
\end{quote}

\begin{quote}
The blacksmith made a new tool for the field.
\end{quote}

Masculine.

\begin{cousinweb}
\textbf{\mr{लोहार}} (\emph{lohār}) is a tadbhava, inherited from Sanskrit \emph{lōhakāra}, \textquotedblleft{}iron-worker\textquotedblright{} (CDIAL 11159), from \emph{lōha}, \textquotedblleft{}iron\textquotedblright{}, and \emph{kāra}, \textquotedblleft{}maker\textquotedblright{}.
\end{cousinweb}

\begin{grammarlens}[title={What you've built}]
One more word for describing what you lived and felt.
\end{grammarlens}

\subsection*{Your turn}
\begin{itemize}
\raggedright
  \item \emph{Say it:} \emph{lohār}
  \item \emph{Say it:} \emph{lohār}, once more
  \item \emph{Say it:} \emph{lohārāne śetāsāṭhī navīn avjār banavlaṁ.} — The blacksmith made a new tool for the field.
  \item \emph{Recall:} say the Marathi for to hammer in, then the Marathi for a potter, then say \emph{lohār} again
\end{itemize}

\begin{tcolorbox}[breakable,colback=teal!4,colframe=teal!35!black,title={Before you move on}]
What does \mr{लोहार} mean? (\textquotedblleft{}A blacksmith\textquotedblright{}.) Say it once more.
\end{tcolorbox}
\section[\texorpdfstring{\mr{सोनार} (sonār)}{सोनार (sonār)}]{\mr{सोनार} (sonār) — a goldsmith}

\label{lesson:MR-C505-sonar}



\begin{quote}
Before the new one: say the Marathi for a potter, then the Marathi for a blacksmith.
\end{quote}

\subsection*{You'll want to know: \mr{सोनार}}
\textbf{\mr{सोनार}} — \emph{sonār} — \textquotedblleft{}a goldsmith\textquotedblright{}.

\begin{quote}
\mr{आईने} \mr{सोनाराकडून} \mr{नवीन} \mr{बांगडी} \mr{घेतली}.
\end{quote}

\begin{quote}
\emph{āīne sonārākaḍūn navīn bāṅgḍī ghetlī.}
\end{quote}

\begin{quote}
Mother bought a new bangle from the goldsmith.
\end{quote}

Masculine.

From a person: -\mr{कडून}, as in \mr{सोनाराकडून}.

\begin{cousinweb}
\textbf{\mr{सोनार}} (\emph{sonār}) is a tadbhava, inherited from Sanskrit \emph{suvarṇakāra}, \textquotedblleft{}goldsmith\textquotedblright{} (CDIAL 13520), from \emph{suvarṇa}, \textquotedblleft{}gold\textquotedblright{}, and \emph{kāra}, \textquotedblleft{}maker\textquotedblright{}.
\end{cousinweb}

\begin{grammarlens}[title={What you've built}]
One more word for describing what you lived and felt.
\end{grammarlens}

\subsection*{Your turn}
\begin{itemize}
\raggedright
  \item \emph{Say it:} \emph{sonār}
  \item \emph{Say it:} \emph{sonār}, once more
  \item \emph{Say it:} \emph{āīne sonārākaḍūn navīn bāṅgḍī ghetlī.} — Mother bought a new bangle from the goldsmith.
  \item \emph{Recall:} say the Marathi for a potter, then the Marathi for a blacksmith, then say \emph{sonār} again
\end{itemize}

\begin{tcolorbox}[breakable,colback=teal!4,colframe=teal!35!black,title={Before you move on}]
What does \mr{सोनार} mean? (\textquotedblleft{}A goldsmith\textquotedblright{}.) Say it once more.
\end{tcolorbox}
\section[\texorpdfstring{\mr{विणकर} (viṇkar)}{विणकर (viṇkar)}]{\mr{विणकर} (viṇkar) — a weaver}

\label{lesson:MR-C505-vinkar}



\begin{quote}
Before the new one: say the Marathi for a blacksmith, then the Marathi for a goldsmith.
\end{quote}

\subsection*{You'll want to know: \mr{विणकर}}
\textbf{\mr{विणकर}} — \emph{viṇkar} — \textquotedblleft{}a weaver\textquotedblright{}.

\begin{quote}
\mr{गावातला} \mr{विणकर} \mr{रेशमी} \mr{साड्या} \mr{विणतो}.
\end{quote}

\begin{quote}
\emph{gāvātlā viṇkar reśmī sāḍyā viṇto.}
\end{quote}

\begin{quote}
The village weaver weaves silk saris.
\end{quote}

Masculine.

Built on \mr{विणणे}, to weave, with -\mr{कर}, one who does.

\begin{cousinweb}
\textbf{\mr{विणकर}} (\emph{viṇkar}) is a Marathi derivation built on \mr{विणणे}, \textquotedblleft{}to weave\textquotedblright{}, with the agent ending -\mr{कर}, \textquotedblleft{}one who does\textquotedblright{}, from Sanskrit \emph{kara} (Monier-Williams).
\end{cousinweb}

\begin{grammarlens}[title={What you've built}]
One more word for describing what you lived and felt.
\end{grammarlens}

\subsection*{Your turn}
\begin{itemize}
\raggedright
  \item \emph{Say it:} \emph{viṇkar}
  \item \emph{Say it:} \emph{viṇkar}, once more
  \item \emph{Say it:} \emph{gāvātlā viṇkar reśmī sāḍyā viṇto.} — The village weaver weaves silk saris.
  \item \emph{Recall:} say the Marathi for a blacksmith, then the Marathi for a goldsmith, then say \emph{viṇkar} again
\end{itemize}

\begin{tcolorbox}[breakable,colback=teal!4,colframe=teal!35!black,title={Before you move on}]
What does \mr{विणकर} mean? (\textquotedblleft{}A weaver\textquotedblright{}.) Say it once more.
\end{tcolorbox}
\section[\texorpdfstring{\mr{थाटणे} (thāṭṇe)}{थाटणे (thāṭṇe)}]{\mr{थाटणे} (thāṭṇe) — to set up; to lay out with care}

\label{lesson:MR-C505-thatne}



\begin{quote}
Before the new one: say the Marathi for a goldsmith, then the Marathi for a weaver.
\end{quote}

\subsection*{You'll want to know: \mr{थाटणे}}
\textbf{\mr{थाटणे}} — \emph{thāṭṇe} — \textquotedblleft{}to set up; to lay out with care\textquotedblright{}.

\begin{quote}
\mr{विणकराने} \mr{बाजारात} \mr{नवीन} \mr{दुकान} \mr{थाटलं}.
\end{quote}

\begin{quote}
\emph{viṇkarāne bājārāt navīn dukān thāṭlaṁ.}
\end{quote}

\begin{quote}
The weaver set up a new shop in the market.
\end{quote}

Past, with -\mr{ने} (he …): \mr{त्याने} \mr{दुकान} \mr{थाटलं}, \emph{tyāne dukān thāṭlaṁ}.

\mr{थाटणे} is said of a shop or a home set up with some show; \mr{थाट} is the show itself.

\begin{cousinweb}
\textbf{\mr{थाटणे}} (\emph{thāṭṇe}) is built on \mr{थाट}, \textquotedblleft{}pomp, a fine arrangement\textquotedblright{}; CDIAL sets Hindi \emph{ṭhāṭh}, \textquotedblleft{}pomp\textquotedblright{}, and Marathi \emph{thaṭṇẽ}, \textquotedblleft{}to act pompously\textquotedblright{}, under \emph{*thaṭṭha}, \textquotedblleft{}pomp, ceremony\textquotedblright{} (CDIAL 6090).
\end{cousinweb}

\begin{grammarlens}[title={What you've built}]
One more word for describing what you lived and felt.
\end{grammarlens}

\subsection*{Your turn}
\begin{itemize}
\raggedright
  \item \emph{Say it:} \emph{thāṭṇe}
  \item \emph{Say it:} \emph{thāṭṇe}, once more
  \item \emph{Say it:} \emph{viṇkarāne bājārāt navīn dukān thāṭlaṁ.} — The weaver set up a new shop in the market.
  \item \emph{Recall:} say the Marathi for a goldsmith, then the Marathi for a weaver, then say \emph{thāṭṇe} again
\end{itemize}

\begin{tcolorbox}[breakable,colback=teal!4,colframe=teal!35!black,title={Before you move on}]
What does \mr{थाटणे} mean? (\textquotedblleft{}To set up; to lay out with care\textquotedblright{}.) Say it once more.
\end{tcolorbox}
